Let \(\mathscr F_m=\sigma(\mathcal P_m)\).  Pilot independence makes \(\hat d_m\), \(\hat\alpha_m\), \(N_{E,B}\), and \(N_{O,B}\) measurable with respect to \(\mathscr F_m\) and independent of the main-study observations.  The pilot-oracle theorem gives
\[
 \hat\alpha_m-\alpha_{\hat d_m}^*(P)=o_P(1).
\]
Finiteness of \(\mathcal D\), positivity of the costs, and nondegeneracy of the variance components therefore give an \(a_0>0\) such that
\[
 \Pr\{a_0\le\hat\alpha_m\le1-a_0\}\longrightarrow1. \tag{1}
\]
On this event \(N_{E,B}\asymp B\) and \(N_{O,B}\asymp B\).  Applying the uniform floor-count result at \(d=\hat d_m\) gives
\[
 s_{0,B}^2-\mathcal V_{\hat d_m}^*(P)=o_P(1),\qquad
 s_{0,B}^2:=B\left\{\frac{V_{E,\hat d_m}}{N_{E,B}}
 +\frac{V_{O,\hat d_m}}{N_{O,B}}\right\}. \tag{2}
\]

We next verify that pilot selection does not disturb the fixed-design expansion.  For a fold \(\ell\), write
\[
 e_{h,d,\ell}=\hat h_{d,B}^{(-\ell)}-h_d,
 \qquad e_{q,d,\ell}=\hat q_{d,B}^{(-\ell)}-q_d.
\]
The exact doubly robust score identity leaves, after subtracting the true-score averages, centered cross-fitted evaluation terms and the bilinear conditional-mean remainder.  For each treatment level, the absolute value of that remainder is bounded, using first \(T_d\) and then its adjoint, by a fixed multiple of
\[
 \min\left\{
 \|e_{h,d,\ell}\|_{w,h,d}\|e_{q,d,\ell}\|_{2,O,q,d},
 \|e_{h,d,\ell}\|_{2,O,h,d}\|e_{q,d,\ell}\|_{w,q,d}
 \right\}. \tag{3}
\]
The four uniform rate assumptions and the weak--strong product condition make (3) \(o_P(B^{-1/2})\), uniformly over the finite family and fixed folds.  Cross-fitting centers the remaining empirical nuisance terms.  The fixed-system asymptotic-distribution theorem recorded in \(\mathrm{lem{:}ikmw\text{-}fixed\text{-}design\text{-}limit\), applied to each member of the finite family, consequently yields the uniform expansion
\[
 \hat\tau_{d,B}-\tau
 =\frac1{n_{E,B}}\sum_{i=1}^{n_{E,B}}\phi_{E,d}(O_{E,i})
 +\frac1{n_{O,B}}\sum_{i=1}^{n_{O,B}}\phi_{O,d}(O_{O,i})
 +o_P(B^{-1/2}) \tag{4}
\]
for allocations in \([a_0,1-a_0]\).  Because the family is finite and the pilot is independent, (4) remains valid after substituting \(d=\hat d_m\) and the pilot-measurable counts.

Conditional on \(\mathscr F_m\), the two leading sums are independent sums of mean-zero variables and their variance after multiplication by \(B\) is exactly \(s_{0,B}^2\).  Each score has finite second moment, and both counts diverge on (1).  The ordinary independent-sample central limit theorem therefore gives
\[
 \frac{\sqrt B\left\{
 N_{E,B}^{-1}\sum_{i=1}^{N_{E,B}}\phi_{E,\hat d_m}(O_{E,i})
 +N_{O,B}^{-1}\sum_{i=1}^{N_{O,B}}\phi_{O,\hat d_m}(O_{O,i})\right\}}
 {s_{0,B}}\rightsquigarrow\mathcal N(0,1). \tag{5}
\]
To justify the random index explicitly, for every fixed source and design the standardized iid mean converges along all integer sample sizes.  Hence it converges uniformly over the tail \(k\ge a_0B/c\).  Finiteness of \(\mathcal D\) makes this uniform in \(d\).  Conditional convergence in (5) is therefore uniform over pilot realizations satisfying (1), and integration over \(\mathscr F_m\) proves the unconditional limit.  Population ties cause no mixture problem because every conditional statistic is divided by its selected-design variance.

Finally, on (1),
\[
\begin{aligned}
 |\hat s_B^2-s_{0,B}^2|
 &\le \frac B{N_{E,B}}|\widetilde V_{E,\hat d_m,B}-V_{E,\hat d_m}|
 +\frac B{N_{O,B}}|\widetilde V_{O,\hat d_m,B}-V_{O,\hat d_m}|\\
 &=O_P(r_{\widetilde V,B})=o_P(1). \tag{6}
\end{aligned}
\]
Equations (2) and (6) prove \(\hat s_B^2-\mathcal V_{\hat d_m}^*(P)=o_P(1)\).  Nondegeneracy and positive costs bound every \(\mathcal V_d^*(P)\) away from zero, so \(\hat s_B/s_{0,B}\to_P1\).  Slutsky's theorem applied to (4)--(6) gives
\[
 \frac{\sqrt B(\hat\tau_B-\tau)}{\hat s_B}\rightsquigarrow\mathcal N(0,1).
\]
If \(z_{1-\eta/2}\) is the corresponding standard-normal quantile, convergence at the two continuity points \(\pm z_{1-\eta/2}\) shows that the stated Wald interval has coverage tending to \(1-\eta\).